\documentclass[10pt,a4paper]{amsart}
\usepackage[margin=19mm]{geometry}
\usepackage{amsmath,amssymb,mathtools}
\usepackage[scr=rsfs,cal=euler]{mathalfa}
\usepackage{xfrac}
\usepackage[unicode,hidelinks,bookmarks=false]{hyperref}
\newcommand{\ie}{i.e.\ }
\newcommand{\cf}{cf.\ }
\newcommand{\resp}{respectively\ }
\newcommand{\Subsection}{\subsection}
\providecommand{\nobreakdash}{\nobreak\hbox{-}}
\setlength{\emergencystretch}{2em}
\multlinegap0pt
\usepackage{bm}
\usepackage{bbm}
\usepackage{dsfont}
\usepackage{xspace}
\usepackage{cancel}
\usepackage{mathtools}
\usepackage{amsthm}
\makeatletter
\@ifundefined{lemma}{\newtheorem{lemma}{Lemma}}{}
\@ifundefined{proposition}{\newtheorem{proposition}{Proposition}}{}
\makeatother
\newcommand{\Mod}[1]{\ (\mathrm{mod}\ #1)}
\title[Kneading the Lorenz attractor]{Kneading the Lorenz attractor}
\author{{\L}ukasz Cholewa, Eran Igra}
\date{Source: arXiv:2511.02568v1 (CC BY 4.0)}
\begin{document}
\maketitle

\noindent\emph{Source and licence.} Kneading the Lorenz attractor. Version arXiv:2511.02568v1 (CC BY 4.0), distributed under the Creative Commons Attribution 4.0 International licence, as recorded in the arXiv licence metadata for this exact version and shown on the abstract page https://arxiv.org/abs/2511.02568v1. Full rights notice: licenses/arxiv-2511.02568-CC-BY-4.0.txt.\par\smallskip

\noindent\textit{Editorial scope.} Counted unit: the Proposition whose source label is \texttt{prop:beta\_seq} in the article (source0.tex lines 1246--1264, source label \texttt{prop:beta\_seq}), delivered together with the article's own complete original proof of it (source0.tex lines 1265--1326, one \texttt{proof} environment of 62 lines). No mathematics is written, repaired or re-derived: statement and proof are the article's own text line for line. Required context retained uncounted: eq:rescaling (lines 581--584); lem:renorm (lines 1122--1128). Every definition, display and label of those lines is retained unchanged. Omitted from this package and not redistributed: 1--580, 585--1121, 1129--1245, 1327--1967. Nothing inside a retained excerpt is omitted or altered: every retained body is delivered exactly as the article's lines are written, so no omission receipt of any kind accompanies this package. Citations: the retained excerpts carry no citation command at all, so no bibliography entry of the article is transcribed and no reference is redistributed. Reference adaptation: none was needed and none was made; every reference inside the delivered unit resolves inside it, so no unresolved reference or citation is produced. Printed slips kept literal: no printed word, symbol, hypothesis or number of the counted statement or of its proof was altered. Layout and class: the article's own class is \texttt{amsart} with packages including graphicx, xargs, geometry, fancyhdr, lastpage, extramarks, tcolorbox, xcolor, hyperref, overpic, pict2e, listings, amssymb, amsthm; the wrapper is the lane's pinned \texttt{amsart} editorial wrapper at 10pt on A4, which restates the theorem-like environments the retained text needs and adds the article's own macro definitions that the retained text needs (\texttt{\textbackslash Mod}), with extra wrapper packages (bm, bbm, dsfont, xspace, cancel, mathtools). The delivered numbering is the unit's own. Assets: no retained span contains a figure environment, a graphics-inclusion command, a PDF-page inclusion, a table or a TikZ picture; no image file is copied into this package, the package declares no asset dependency and no image file is redistributed with it. Licence: version arXiv:2511.02568v1 of this article is distributed under the Creative Commons Attribution 4.0 International licence, as recorded in the arXiv licence metadata for this exact version and shown on the abstract page https://arxiv.org/abs/2511.02568v1. A full rights notice is delivered at licenses/arxiv-2511.02568-CC-BY-4.0.txt. Characters: every retained excerpt is pure ASCII, so the delivered engine, XeLaTeX with its default Latin Modern text font, reports no missing character.
\begin{equation}
\label{eq:rescaling}
h_{u,v}\colon[u,v]\to[0,1],\quad h_{u,v}(x)=\frac{x-u}{v-u}.
\end{equation}

\begin{lemma}\label{lem:renorm}
    Let $F\colon[0,1]\to[0,1]$ be an expanding Lorenz map with a primary $n(k)$-cycle. Then the set of periods $P_F$ of the map $F$ is described by the following formula
    $$
    P_F=n\cdot P_G\cup\left\lbrace n \right\rbrace=\left\lbrace n\cdot q\:|\: q\in P_G \right\rbrace\cup\left\lbrace n\right\rbrace,
    $$
    where $P_G$ is the set of periods of the renormalization $G=(F^n,F^n)$.
\end{lemma}

\begin{proposition}
\label{prop:beta_seq}
    Let $\beta\in(\beta_{i+1},\beta_i]$ for some $i\in\mathbb{N}$. Then the following conditions hold.
    \begin{enumerate}
    \item\label{prop:beta_seq1} The map $F_{\beta}$ has $i$ renormalizations 
    $$
    G_{\beta,j}=(F^{2^j}_{\beta},F^{2^j}_{\beta}),\quad\text{for}\quad j=1,\ldots,i,
    $$ 
    that after rescaling to $[0,1]$ are equal to the maps $F_{\beta^{2^j}}$ defined by the parameters $\beta^{2^j}\in(\beta_{i+1-j},\beta_{i-j}]$.     
    \item\label{prop:beta_seq2} The set of periods of $F_{\beta}$ is given by
    $$
    P_{F_{\beta}}=2\cdot P_{F_{\beta^{2}}}\cup\{2\}.
    $$
    Moreover, if $i\geq2$, then
    $$
    P_{F_{\beta}}=2^{j}\cdot P_{F_{\beta^{2^j}}}\cup\left\lbrace2^{j-1},\ldots,2\right\rbrace,\quad\text{for}\quad j=2,\ldots,i.
    $$
    \end{enumerate}
\end{proposition}

\begin{proof}
The proof is by induction on $i\in\mathbb{N}$. First, observe that for $i=1$ the conditions~\eqref{prop:beta_seq1} and \eqref{prop:beta_seq2} follow immediately from Lemma~\ref{lem:renorm}.

Assume that there is $i$ such that the conditions~\eqref{prop:beta_seq1} and \eqref{prop:beta_seq2} holds for every parameter from $(\beta_{i+1},\beta_i]$. Let $\beta\in(\beta_{i+2},\beta_{i+1}]$. Since $(\beta_{i+2},\beta_{i+1}]\subset(1,\sqrt{2}]$, the map $F_{\beta}$ has renormalization $G_{\beta,1}:=G_{\beta}=(F^2_{\beta},F^2_{\beta})$, that is
$$
G_{\beta}(x)=
\begin{cases}F^2 (x),\,&\text{if 
		$x\in\left[u,c\right)$}\\
	F^2 (x),\,&\text{if $x\in (c,v]$}\end{cases}
=
\begin{cases}\beta^2x+\frac{\beta-\beta^2}{2},\,&\text{if 
		$x\in\left[u,c\right)$}\\
	\beta^2x+1-\frac{\beta+\beta^2}{2},\,&\text{if $x\in (c,v]$}\end{cases},
$$
where $[u,v]=[F_{\beta}(0),F_{\beta}(1)]$ and $c=\frac{1}{2}$. Let $h_{u,v}\colon[u,v]\to[0,1]$ be the map defined by~\eqref{eq:rescaling}. Simple calculations yield that 
$$
\begin{aligned}
\left(h_{u,v}\circ G_{\beta}\circ h^{-1}_{u,v}\right)(x)&=
\begin{cases}h_{u,v}\left(\beta^2h_{u,v}^{-1}(x)+\frac{\beta-\beta^2}{2}\right),\,&\text{if 
		$x\in\left[0,c\right)$}\\
	h_{u,v}\left(\beta^2h_{u,v}^{-1}(x)+1-\frac{\beta+\beta^2}{2}\right),\,&\text{if $x\in (c,1]$}\end{cases}\\
    &=
\begin{cases}\beta^2x+1-\frac{\beta^2}{2},\,&\text{if 
		$x\in\left[0,c\right)$}\\
	\beta^2x-\frac{\beta^2}{2},\,&\text{if $x\in (c,1]$}\end{cases}\\
    &=\beta^2 x+1-\frac{\beta^2}{2}\Mod{1}=F_{\beta^2}(x).
\end{aligned}
$$
So the renormalization $G_{\beta,1}$ after rescaling to $[0,1]$ is equal to the map $F_{\beta^2}$ and $\beta^2\in(\beta_{i+1},\beta_i]$.

By the induction assumption the map $F_{\beta^2}$ has $i$ renormalizations 
$$
    G_{\beta^2,j}=(F^{2^j}_{\beta^2},F^{2^j}_{\beta^2}),\quad\text{for}\quad j=1,\ldots,i,
$$ 
that (after rescaling) are equal to the maps $F_{(\beta^2)^{2^j}}=F_{\beta^{2^{j+1}}}$, where $\beta^{2^{j+1}}\in(\beta_{i+1-j},\beta_{i-j}]$. Fix $j\in\{1,\ldots,i\}$ and denote by $[u_j,v_j]$ the renormalization interval of $G_{\beta^2,j}$. Then, for every $x\in[u_j,v_j]$ we have
$$
G_{\beta^2,j}(x)=F^{2^j}_{\beta^2}(x)=\left(h_{u,v}\circ G_{\beta}\circ h^{-1}_{u,v}\right)^{2^j}(x)=\left(h_{u,v}\circ G_{\beta}^{2^j}\circ h^{-1}_{u,v}\right)(x)=\left(h_{u,v}\circ F_{\beta}^{2^{j+1}}\circ h^{-1}_{u,v}\right)(x).
$$
Since $h_{u_j,v_j}\circ G_{\beta^2,j}\circ h^{-1}_{u_j,v_j}=F_{\beta^{2^{j+1}}}$, we obtain
$$
F_{\beta}^{2^{j+1}}(x)=\left(h_{u,v}^{-1}\circ G_{\beta^2,j}\circ h_{u,v}\right)(x)=\left[(h_{u_j,v_j}\circ h_{u,v})^{-1}\circ F_{\beta^{2^{j+1}}}\circ(h_{u_j,v_j}\circ h_{u,v})\right](x)
$$
for $x\in[a,b]:=h_{u,v}^{-1}([u_j,v_j])=[h_{u,v}^{-1}(u_j),h_{u,v}^{-1}(v_j)]$. Next, observe that $h_{u_j,v_j}\circ h_{u,v}\colon[a,b]\to[0,1]$ and
$$
\left(h_{u_j,v_j}\circ h_{u,v}\right)(x)=\frac{x-(u_j(v-u)+u)}{(v_j-u_j)(v-u)}=\frac{x-h_{u,v}^{-1}(u_j)}{h_{u,v}^{-1}(v_j)-h_{u,v}^{-1}(u_j)}=\frac{x-a}{b-a}=h_{a,b}(x).
$$
Therefore $h_{a,b}\circ F_{\beta}^{2^{j+1}}\circ h^{-1}_{a,b}=F_{\beta^{2^{j+1}}}$, so the map $F_{\beta}^{2^{j+1}}\colon[a,b]\to[a,b]$ is an expanding Lorenz map on $[a,b]$. Hence $G_{\beta,j+1}:=(F_{\beta}^{2^{j+1}},F_{\beta}^{2^{j+1}})$ is a renormalization of $F_{\beta}$ for each $j=1,\ldots,i$.

Finally, let us note that by Lemma~\ref{lem:renorm} we have
$$
P_{F_{\beta}}=2\cdot P_{F_{\beta^{2}}}\cup\{2\}.
$$
Using again the induction assumption, we obtain
$$
P_{F_{\beta}}=2\cdot\left(2\cdot P_{F_{(\beta^2)^{2}}}\cup\{2\}\right)\cup\{2\}=2^2\cdot P_{F_{\beta^{2^2}}}\cup\{2^2,2\}
$$
and
$$
P_{F_{\beta}}=2\cdot\left(2^{j}\cdot P_{F_{(\beta^2)^{2^j}}}\cup\left\lbrace2^{j-1},\ldots,2\right\rbrace\right)\cup\{2\}=2^{j+1}\cdot P_{F_{\beta^{2^{j+1}}}}\cup\left\lbrace2^{j},\ldots,2\right\rbrace,
$$
for $j=2,\ldots,i$. The proof is completed.
\end{proof}

\end{document}
